\section{Introduction} \label{sec:introduction}

With the rapid development of the Internet of Things, massive real-time data
streams have become ubiquitous in a variety of modern applications such as
traffic management and intelligent parking~\cite{Babcock2002, Babu2001}.
However, the collection and analysis of sensitive user data, such as real-time
vehicle locations in a city and user purchasing history on e-commerce websites,
pose significant privacy risks~\cite{Al2022, Ray2020, Malhotra2021}.
To protect user privacy, differential privacy (DP)~\cite{Dwork2006} has emerged
as a cornerstone, offering rigorous guarantees by injecting calibrated noise into
user data.
While centralized differential privacy (CDP)~\cite{Bolot2013} relies on a
trusted server to collect and perturb data, local differential privacy
(LDP)~\cite{Wang2017l, Wang2018p, Ren2018, Zhang2018, Duchi2013, Cormode2018,
  Wang2019l, Murakami2019, Ye2019, Qin2016, Qin2017} eliminates this requirement
by applying perturbation directly on the user side.
This advantage has led to the deployment of LDP in large-scale systems,
including Google's RAPPOR~\cite{Erlingsson2014}, Apple’s iOS
analytics~\cite{Team2017}, and Microsoft’s telemetry services~\cite{Ding2017}.

Researchers have proposed many privacy-preserving methods for streaming data,
most of which focus on user-level DP for finite streams~\cite{Fan2013a,
  Fan2013, Fan2012} or event-level DP for infinite streams~\cite{Dwork2010,
  Bolot2013}.
However, real-world applications usually run for long periods, producing
infinite streams.
Providing event-level privacy on such streams cannot satisfy the requirements for
protecting arbitrary events.
In contrast, providing user-level privacy requires adding infinite
perturbations, which ultimately reduces the utility of the data.
To address these issues, Kellaris et al.~\cite{Kellaris2014} proposed $w$-event
DP to protect any sequence of events that occur in any $w$ consecutive
timestamps (i.e., a sliding window of size $w$), achieving a balance between
privacy and utility for infinite streams.

However, existing $w$-event DP studies~\cite{Ren2020,Wang2020,Wang2019,Wang2016}
are primarily designed under the CDP setting.
As an emerging solution for infinite data streams, $w$-event LDP~\cite{Ren2022}
shows strong potential but remains underexplored.
There are two key limitations that remain understudied and hinder its practical
adoption.
First, existing mechanisms fail to achieve efficient privacy budget utilization.
While dynamic budget allocation is theoretically possible, empirical studies
reveal that these strategies often perform worse than simple baselines, such as
uniform allocation or random sampling~\cite{Schaler2024}.
Second, $w$-event LDP is severely limited in the types of analytical tasks it
supports.
While event-level and user-level DP methods have matured to support diverse
analytical tasks including counting, range queries, and event
monitoring~\cite{Chen2017, Cao2017, Chan2012}, existing $w$-event methods remain
limited to fixed-granularity statistical releases, particularly frequency
histograms~\cite{Ren2022,Li2025}.
This single-focus output makes it impossible to perform \emph{multi-granularity}
analysis on the same continuous stream.
This rigidity is a critical barrier in real-world applications such as
intelligent vehicle systems~\cite{sharif2017internet},
where operators require
insights at different \emph{spatial} resolutions.
For example, they require fine-grained insights, such as the specific vehicle
count at a single intersection to manage traffic lights.
Simultaneously, they must analyze coarse-grained data, such as the total traffic
volume along an entire arterial road to identify bottlenecks.

To overcome these limitations and design a $w$-event LDP framework that supports
concurrent multi-task processing, we identify four fundamental challenges to be
addressed:
\begin{itemize}[nosep]
\item \textbf{Spatial Imbalance in Data Distribution.}
  Real-world streaming data often exhibits significant spatial imbalance within
  individual timestamps.
  For instance, in intelligent transportation systems, traffic on arterial roads
  may far exceed that on side roads.
  Existing methods that uniformly partition the domain and apply the same noise scale to all regions can bury signals from sparse regions beneath the noise magnitude, obscuring underlying patterns and impairing downstream analytics.

%动态层次结构
\item \textbf{Temporal Variation in Data Distribution.}
  The variation in stream distributions over time poses a significant challenge for
  allocating privacy budgets that can adaptively adjust to the degree of
  variation across the sliding window.
  Existing methods, often focusing solely on local information, fail to capture these dynamics, resulting in poor utility.

\item \textbf{Privacy Budget Sharing across Multiple Tasks.}
  Streaming systems are often required to support multiple query tasks
  concurrently.
  As the number of concurrent tasks increases, the limited privacy budget must
  be divided among more tasks, causing a rapid decline in the budget allocated
  to each task and significantly degrading overall accuracy.

\item \textbf{Low-Latency Constraints in Real-Time Processing.}
  Real-time streaming applications impose strict latency requirements.
  Even if a more accurate mechanism exists, it may be impractical if it cannot
  respond within tight deadlines.
  This necessitates a trade-off among privacy protection, estimation accuracy,
  and processing latency.
\end{itemize}

To address the above challenges, we propose MTSP-LDP, a new $w$-event LDP
framework for infinite streams.
Our main contributions are summarized as follows.
\begin{itemize}[nosep]
\item {To the best of our knowledge, this paper is the first to address the
  challenge of supporting multi-task, multi-granularity analytical tasks under
  $w$-event LDP.
  Specifically, we introduce the \emph{Budget-Free Multi-Task Processing}
  mechanism, which enables concurrent support for different types of queries
  over infinite data streams.
  Notably, this mechanism operates without requiring modifications to other
  modules or the perturbation data reported by users.
  Crucially, it incurs no additional privacy cost.}

\item {Our framework's high utility is achieved through two key technical
  innovations.
  First, the \emph{Optimal Privacy Budget Allocation} mechanism adaptively
  allocates privacy budgets based on data variation across each sliding window
  of size $w$, ensuring that timestamps with greater variation receive more
  privacy budget.
  Second, the \emph{Private Adaptive Tree Publication} mechanism improves accuracy by constructing a tree tailored to the underlying data distribution, departing from traditional fixed-partition histograms. It then applies cross-timestamp grouping and smoothing to further enhance estimation accuracy. Crucially, our mechanism resolves the prohibitive level-by-level interaction latency associated with traditional adaptive trees originally designed for static settings, thus enabling data-adaptive structures for streaming data.}

\item We conduct extensive experiments on four real-world datasets.
  Experimental results demonstrate that MTSP-LDP consistently outperforms
  state-of-the-art approaches designed for these specific tasks, achieving high
  utility under strict privacy guarantees.
\end{itemize}

% This paper is organized as follows.
% Section \ref{sec:Preliminaries} reviews the relevant background knowledge.
% In Section \ref{sec:Problem}, we formulate the data model and query tasks on
% streams.
% In Section \ref{sec:method}, we describe the framework of MTSP-LDP and
% explains how it supports multiple query tasks.
% Section \ref{sec:Evaluation} presents comprehensive experimental results.
% Section \ref{sec:related} discusses related work.
% Section \ref{sec:conclusions} is the conclusion of this paper.
